\subsection{可度量化向量空间中的连续线性函数}

定理 1 (Banach). --- 设 E 与 F 是非离散赋值除环 K 上的两个可度量化向量空间，并设 u 是 E 到 F 中的一个连续线性映射。假设 E 是完备的。则下列条件等价：

\begin{enumerate}
\def\labelenumi{(\roman{enumi})}
\item
  \(u\) 是一个严格满射态射。
\item
  F 是完备的且 \(u\) 是满射。
\item
  \(u\) 的像在 \(\mathbf{F}\) 中不是贫集（GT, IX, § 5.2）。
\item
  对 0 的每个邻域 \(\mathbf{V}\)（在 \(\mathbf{E}\) 中），集合 \(\overline{u(\mathbf{V})}\) 是 0 在 \(\mathbf{F}\) 中的一个邻域。首先 (i) 蕴含 (ii)，因为设 \(u\) 是一个严格满射态射，\(\mathbf{N}\) 是 \(u\) 的核。则 \(u\) 诱导 \(\mathrm{E/N}\) 到 \(\mathbf{F}\) 上的一个同构。但 \(\mathbf{E}\) 是可度量化且完备的，故 \(\mathrm{E/N}\) 是完备的（GT, IX, § 3.1, prop. 4），因此 \(\mathbf{F}\) 是完备的。
\end{enumerate}

其次 (ii) 蕴含 (iii)。设 F 是完备的且 u 是满射。u 的像恰好是 F，因而由 Baire 定理（GT, IX, § 5.3）它在 F 中不是贫集。

下述引理表明 (iii) 蕴含 (iv)。

引理 1. --- 设 E 与 F 是非离散赋值除环 K 上的两个拓扑向量空间，并设 u 是 E 到 F 中的一个连续线性映射，且 E 的像不是贫集。则对 E 中 0 的每个邻域 V，集合 \(\overline{u(V)}\) 是 F 中 0 的一个邻域。

设 W 是 E 中 0 的一个均衡邻域，满足 \(W + W \subset V\)（I, § 1.5, prop. 4）。设 \(\alpha\) 是 K 中满足 \(|\alpha| > 1\) 的一个元素；则 E 是诸集合 \(\alpha^{n}W\)（其中 n 取遍 N）的并；事实上，对每个 \(x \in E\)，存在 \(\beta \in K\) 使 \(x \in \beta W\)（I, 第 7 页, prop. 4），且存在整数 \(n \geqslant 0\) 使 \(|\beta| < |\alpha|^{n}\)，于是由 W 是均衡的，得 \(x \in \alpha^{n}W\)。因此，\(u(E)\) 是集合序列 \(u(\alpha^{n}W) = \alpha^{n}u(W)\) 的并；由于 \(u(E)\) 在 F 中不是贫集，诸集合 \(\alpha^{n}\overline{u(W)}\) 中至少有一个具有内点（GT, IX, § 5.3, def. 2），从而 \(\overline{u(W)}\) 具有内点。

设 \(y_0\) 是 \(\overline{u(\mathbf{W})}\) 的一个内点；由于 \(-u(\mathbf{W}) = u(\mathbf{W})\)，从而 \(-\overline{u(\mathbf{W})} = \overline{u(\mathbf{W})}\)，由此得出 \(0 = y_0 + (-y_0)\) 是 \(\overline{u(\mathbf{W})} + \overline{u(\mathbf{W})}\) 的一个内点。由于向量加法是 \(F \times F\) 到 \(F\) 中的连续映射，集合 \(\overline{u(\mathbf{W})} + \overline{u(\mathbf{W})}\) 包含在集合

\[
u (\mathrm{W}) + u (\mathrm{W}) = u (\mathrm{W} + \mathrm{W}) \subset u (\mathrm{V});
\]

的闭包之中；因此 \(\overline{u(\mathrm{V})}\) 是 F 中 0 的一个邻域。

在证明 (iv) 蕴含 (i) 之前，我们先证明下述引理，其中我们约定：在所有度量空间中，\(\mathbf{B}_{r}(x)\) 表示以 x 为中心、以 r 为半径的闭球。

引理 2. --- 设 E 与 F 是两个度量空间，且假设 E 还是完备的。设 u 是 E 到 F 中的一个线性映射，具有下述性质：无论数 r \textgreater{} 0 如何，都存在一个数 \(\rho(r) > 0\)，使得对每个 \(x \in E\)，我们有

\[
\mathrm{B} _ {\rho (r)} (u (x)) \subset \overline {{u (\mathrm{B} _ {r} (x))}}.
\]

在这些条件下，对每个 \(a > r\)，像 \(u(\mathbf{B}_a(x))\) 包含球 \(\mathbf{B}_{\rho(r)}(u(x))\)。

设 \((r_n)\) 是一个无穷的数列 \(>0\)，满足 \(r_1 = r\) 且 \(a = \sum_{n=1}^{\infty} r_n\)。对每个指标 \(n\)，存在一个数 \(\rho_n > 0\)（其中 \(\rho_1 = \rho(r)\)）使得

\[
\mathrm{B} _ {\rho_ {n}} (u (x)) \subset \overline {{u (\mathrm{B} _ {r _ {n}} (x))}}
\]

对每个 \(x \in E\)；我们可以而且将假设 \(\lim \rho_{n} = 0\)。

设 \(x_0\) 是 \(E\) 的一个点，\(y\) 是 \(B_{\rho(r)}(u(x_0))\) 的一个点。我们将证明 \(y\) 属于 \(u(B_a(x_0))\)。

为此，对 E 的点归纳地定义一个序列 \((x_{n})_{n>0}\)，使得对每个 \(n \geqslant 1\)，有 \(x_{n} \in \mathbf{B}_{r_{n}}(x_{n-1})\) 且 \(u(x_{n}) \in \mathbf{B}_{\rho_{n+1}}(y)\)。若诸 \(x_{i}\) 已对 \(0 \leqslant i \leqslant n - 1\) 按满足这些关系的方式定义好，则我们有 \(y \in \mathbf{B}_{\rho_{n}}(u(x_{n-1}))\)；由于

\[
\mathrm{B} _ {\rho_ {n}} (u (x _ {n - 1})) \subset \overline {{u (\mathrm{B} _ {r _ {n}} (x _ {n - 1}))}},
\]

则存在一个点 \(x_{n} \in \mathbf{B}_{r_n}(x_{n-1})\)，它的像 \(u(x_{n})\) 属于邻域 \(\mathbf{B}_{\rho_{n+1}}(y)\)（\(y\) 的邻域），这就确立了序列 \((x_{n})\) 的存在性。

由于 \(x_{n}\) 与 \(x_{n+p}\) 的距离小于 \(r_{n+1} + r_{n+2} + \cdots + r_{n+p}\)（当 n 充分大时它可任意小），序列 \((x_{n})\) 是 E 中的 Cauchy 序列。由于 E 是完备的，序列 \((x_{n})\) 收敛于 E 的一个点 x。\(x_{0}\) 与 x 的距离小于 \(\sum_{n=1}^{\infty} r_{n} = a\)，从而 \(x \in \mathbf{B}_{a}(x_{0})\)。但 u 是连续的，故序列 \(u(x_{n})\) 收敛于 \(u(x)\)；又有 \(u(x_{n}) \in \mathbf{B}_{\rho_{n+1}}(y)\)，因此 \(y = u(x)\)，引理得证。

我们回到定理并证明 (iv) 蕴含 (i)。设 \(u\) 满足条件 (iv)。对空间 E 与 F 中的每一个，取一个在平移下不变且定义其拓扑的距离（I, 第 16 页）。根据假设，集合 \(u(\mathbf{B}_r(0))\) 对每个 \(r > 0\) 都是 F 中 0 的邻域，从而存在一个数 \(\rho(r) > 0\) 使 \(\mathbf{B}_{\rho(r)}(0) \subset u(\mathbf{B}_r(0))\)。通过平移我们断言：\(\mathbf{B}_{\rho(r)}(u(x)) \subset u(\mathbf{B}_r(x))\) 对每个 \(r > 0\) 及每个 \(x \in E\) 成立。由引理 2，对每一对正实数 \((a, r)\)（\(a > r > 0\)），我们有 \(\mathbf{B}_{\rho(r)}(0) \subset u(\mathbf{B}_a(0))\)；因此 \(u\) 是 E 到 F 上的一个严格态射。我们已证明 (iv) 蕴含 (i)，定理的证明至此完成。

推论 1. --- 若 E 与 F 是非离散赋值除环上的两个完备可度量化向量空间，则 E 到 F 上的每个双射连续线性映射都是同构。

特别地，若 E 与 F 是两个完备赋范空间，则存在一个数 \(a > 0\)，使得 \(\| u(x)\| \geqslant a.\| x\|\) 对每个 \(x\in \mathbf{E}\) 成立。

推论 2. --- 设 E 是非离散赋值除环上的一个向量空间，\(T_{1}\) 与 \(T_{2}\) 是 E 上与其向量空间结构相容的两个拓扑，且对其中每个拓扑 E 都是可度量化且完备的。则若 \(T_{1}\) 与 \(T_{2}\) 可比较，它们就恒同。

推论 3. --- 设 E 与 F 是非离散赋值除环上的两个完备可度量化向量空间。为使 E 到 F 中的一个连续线性映射 u 是严格态射，必须且只需 \(u(\mathrm{E})\) 在 F 中是闭的。

该条件是必要的，因为若 u 是严格态射，则像 \(u(\mathrm{E})\) 同构于商 \(\mathrm{E}/u^{-1}(0)\)，因而是完备的（I, 第 17 页），从而在 F 中是闭的。该条件是充分的，因为若 \(u(\mathrm{E})\) 在 F 中是闭的，则 \(u(\mathrm{E})\) 必是一个完备可度量化向量空间，从而由定理 1，u 是 E 到 \(u(\mathrm{E})\) 上的一个严格态射。

推论 4. --- 设 E 是非离散赋值除环上的一个完备可度量化向量空间。若 M 与 N 是 E 的两个闭向量子空间，且它们在 E 中（代数）互补，则 E 是 M 与 N 的拓扑直和。

因为 \(M \times N\) 是完备可度量化向量空间，且映射 \((y, z) \mapsto y + z\)（\(M \times N\) 到 E 上的映射）连续且双射，从而是同构（cor. 1）。

推论 5 (闭图像定理). --- 设 E 与 F 是非离散赋值除环上的两个完备可度量化向量空间。为使 E 到 F 中的一个线性映射 u 连续，必须且只需它的图在乘积空间 \(E \times F\) 中是闭的。

该条件是必要的，因为到 Hausdorff 空间中的连续映射的图是闭的（GT, I, § 8.1, cor. 2）。为看出它是充分的，注意它蕴含：u 的图 G（它是完备可度量化空间 \(E \times F\) 的一个闭向量子空间）本身是可度量化的且完备的。G 到 E 上的投影 \(z \mapsto \operatorname{pr}_{1}(z)\) 是一个双射的连续线性映射，从而是同构（cor. 1）；由于它的逆映射是 \(x \mapsto (x, u(x))\)，由此得出 \(u\) 在 E 中连续。

我们可以把这个推论表述成如下形式：\(u\) 是连续的，如果下述情形成立：若 E 的点的序列 \((x_{n})\) 既收敛于 0，又使序列 \((u(x_{n}))\) 收敛于 \(y\)，则必然有 \(y = 0\)。

例. --- 设 E 是定义在 I = \([0, 1]\) 上的实值函数所成空间的一个向量子空间；设 \(\|f\|\) 是 E 上的一个范数，在此范数下 E 是完备的，且使得 E 的拓扑比简单收敛拓扑更细。进一步假设 E 包含在 I 上无穷次可微的函数所成的集合 \(\mathcal{C}^{\infty}(\mathrm{I})\)；我们将证明存在一个整数 \(k \geqslant 0\)，使 E 包含在 I 中具有连续的 k 阶导数的所有函数所成的集合 \(\mathcal{C}^{k}(\mathrm{I})\)。

对每一对整数 \(m > 0\), \(n \geqslant 0\)，设 \(V_{mn}\) 是由那些 \(f \in \mathcal{C}^{\infty}(\mathrm{I})\)（满足 \(|f^{(h)}(x)| \leqslant 1 / m\)，其中 \(0 \leqslant h \leqslant n\) 且对每个 \(x \in \mathrm{I}\)）所成的集合。诸 \(V_{m,n}\) 构成与 \(\mathcal{C}^{\infty}(\mathrm{I})\) 的向量空间结构相容的某个可度量化拓扑的 0 的一个基本邻域系，此外 \(\mathcal{C}^{\infty}(\mathrm{I})\) 在此拓扑下是完备的（FVR, II, 第 2 页, th. 1）。设 \(u\) 是 \(\mathcal{C}^{\infty}(\mathrm{I})\) 到 E 中的典范映射；我们证明 \(u\) 连续。由上面的 cor. 5，只需证明：若一个序列 \((f_n)\) 在 \(\mathcal{C}^{\infty}(\mathrm{I})\) 中收敛于 0 且在 E 中收敛于一个极限 \(f\)，则必然 \(f = 0\)。但这是直接的，因为根据假设，\(f\) 是 \((f_n)\) 的简单收敛极限。因此存在一个整数 \(k \geqslant 0\) 和一个数 \(a > 0\)，使得关系

\[
p_{k}(f) = \sup_{\substack{x\in \mathbf{I}\\ 0\leqslant h\leqslant k}}\big|f^{(h)}(x)\big|\leqslant a
\]

蕴含 \(\| f\| \leqslant 1\)，对每个 \(f\in \mathcal{C}^{\infty}(\mathrm{I})\)

但 \(p_k\) 是空间 \(\mathcal{C}^k(\mathrm{I})\) 上的一个范数，且对此范数 \(\mathcal{C}^\infty(\mathrm{I})\) 是在 \(\mathcal{C}^k(\mathrm{I})\) 中处处稠密的一个子空间（多项式的集合已在 \(\mathcal{C}^k(\mathrm{I})\) 中处处稠密，这是 Weierstrass-Stone 定理的一个直接推论）。由前所述，\(\mathcal{C}^\infty(\mathrm{I})\)（带范数 \(p_k\)）到 \(\mathrm{E}\) 中的恒等映射是连续的，因而它可以连续地延拓到整个空间 \(\mathcal{C}^k(\mathrm{I})\)（因为 \(\mathrm{E}\) 是完备的）。这就证明了我们的断言。

命题 1. --- 设 E, F 是非离散赋值除环 K 上的两个拓扑向量空间。我们假设：

\begin{enumerate}
\def\labelenumi{\arabic{enumi})}
\item
  E 是可度量化且完备的。
\item
  存在 K 上的完备可度量化向量空间的一个序列 \((\mathbf{F}_{n})\)，且对每个 n，存在一个单射连续线性映射 \(v_{n}\)（\(F_{n}\) 到 F 中的映射），使得 F 是诸子空间 \(v_{n}(\mathbf{F}_{n})\) 的并。
\end{enumerate}

再设 u 是 E 到 F 中的一个线性映射。若 u 的图在 \(E \times F\) 中是闭的，则存在一个整数 n 和一个连续线性映射 \(u_{n}\)（E 到 \(F_{n}\) 中的映射），使得 \(u = v_{n} \circ u_{n}\)（这蕴含 u 连续且 \(u(\mathrm{E}) \subset v_{n}(\mathrm{F}_{n})\)）。

设 G 是 u 在 \(E \times F\) 中的图。对每个 n，我们考虑连续线性映射 \(w_{n}:(x,y)\mapsto(x,v_{n}(y))\)（\(E \times F_{n}\) 到 \(E \times F\) 中的映射）；由于 G 是闭的，集合 \(w_{n}^{-1}(\mathbf{G})=\mathbf{G}_{n}\) 是 \(E \times F_{n}\) 的一个闭向量子空间；若 \(p_{n}\) 是 \(G_{n}\) 上的限制映射，即第一投影 \(pr_{1}\) 的限制，则我们有 \(p_{n}(\mathbf{G}_{n})=u^{-1}(v_{n}(\mathbf{F}_{n}))\)。由于 \(p_{n}\) 连续且 \(G_{n}\) 完备（因为 \(G_{n}\) 在完备空间 \(E \times F_{n}\) 中是闭的），由定理 1，\(p_{n}(\mathbf{G}_{n})\) 或者在 E 中是贫集，或者是整个 E。但根据假设，E 是诸 \(p_{n}(\mathbf{G}_{n})\) 的并，且由于 E 是完备的，由 Baire 定理（GT, IX, § 5.3, th. 1），诸 \(p_{n}(\mathbf{G}_{n})\) 不可能全是 E 中的贫集。因此存在一个整数 n 使 \(p_{n}(\mathbf{G}_{n}) = \mathbf{E}\)，换言之 \(u(\mathbf{E}) \subset v_{n}(\mathbf{F}_{n})\)。此外，由于 \(v_{n}\) 是单射，\(G_{n}\) 是某个线性映射 \(u_{n}\)（E 到 \(F_{n}\) 中的映射）的图，且由闭图像定理（I, 第 19 页, cor. 5）\(u_{n}\) 连续；于是由定义可得 \(u = v_{n} \circ u_{n}\)。

